\documentclass[11pt,a4paper]{article}
\usepackage[margin=26mm]{geometry}
\usepackage{amsmath,amssymb,amsthm}
\usepackage[T1]{fontenc}
\usepackage{lmodern}
\usepackage{hyperref}
\newtheorem{proposition}{Proposition}
\title{Infinitely many highly transitive complete designs\\\large Counterexample to TLMC 00000008411}
\author{ziangni-sys}
\date{4 October 2026}
\begin{document}
\maketitle
\begin{abstract}
For every $v\geq5$, the complements of single points in a $v$-point
set form a $3$-$(v,v-1,v-3)$ design. Every point permutation is an
automorphism, and the automorphism group is $3$-transitive. Different
values of $v$ give nonisomorphic designs. This supplies infinitely many
isomorphism classes with the fixed strength $t=3$.
\end{abstract}

\section{Scope of the assertion}
The source defines a $t$-$(v,k,\lambda)$ design by the condition that
each $t$-element subset is contained in exactly $\lambda$ blocks. It
then asserts that for fixed $t\geq3$ only finitely many isomorphism
classes have a $t$-transitive automorphism group.

The following family consists of \emph{complete designs}, often called
trivial designs in classification terminology. The stated definition
does not exclude complete designs. Our parameters nevertheless satisfy
$3<k<v$. We make no claim about a different classification problem that
adds a noncomplete-design hypothesis. Only $t$ is fixed in the source;
the parameters $v,k,\lambda$ are allowed to vary.

\section{An explicit infinite family}
For each $v\geq5$, let $X_v$ be a set of $v$ points and put
\[
 \mathcal B_v=\{X_v\setminus\{x\}:x\in X_v\}.
\]
Distinct omitted points give distinct blocks, so this is a simple
block system with exactly $v$ blocks, each of cardinality $v-1$.
It is precisely the family of all $(v-1)$-element subsets of $X_v$.

\begin{proposition}
$(X_v,\mathcal B_v)$ is a $3$-$(v,v-1,v-3)$ design, and its
automorphism group is $3$-transitive.
\end{proposition}
\begin{proof}
Fix any three-element subset $T\subseteq X_v$. The block
$X_v\setminus\{x\}$ contains $T$ exactly when $x\notin T$.
There are $v-3$ such points, and distinct points give distinct blocks.
This proves the design property with $\lambda=v-3>0$.

For every permutation $\sigma$ of $X_v$,
\[
 \sigma(X_v\setminus\{x\})=X_v\setminus\{\sigma(x)\}.
\]
Thus every permutation preserves the block family and is an
automorphism. Given any two ordered triples of distinct points, send
the first triple to the second in the prescribed order. This bijection
between their three-element ranges extends to a permutation of $X_v$,
because the two complements have equal cardinality $v-3$.
The resulting automorphism sends the triples to one another, proving
$3$-transitivity.
\end{proof}

\newpage
\begin{proposition}
These designs have infinitely many distinct isomorphism classes.
\end{proposition}
\begin{proof}
An isomorphism of block designs includes a bijection of their point
sets. Consequently the designs for $v$ and $w$ can be isomorphic only
if $v=w$. Taking $v=n+5$ for every natural number $n$ therefore gives
an infinite sequence of pairwise nonisomorphic designs. The value
$t=3$ is fixed throughout, so this contradicts the claimed finiteness.
\end{proof}

The finiteness clause already fails, so the subsequent classification
list cannot establish the conjecture as stated.

\section{Formal correspondence}
\raggedright
The Lean project defines the actual point set as \texttt{Fin (n+5)}.
The block omitting $i$ is \texttt{Finset.univ.erase i}, and the block
family is its finite image over all points. The theorem
\texttt{block\_injective} rules out duplicate blocks, while
\texttt{block\_card} and \texttt{number\_of\_blocks} verify the sizes.

The definition \texttt{IsThreeDesign} requires every block to have
$n+4$ points and quantifies over every actual three-element subset,
counting the blocks containing it. The theorem
\texttt{incidence\_count} identifies this filtered family with the
omitted points outside that subset and computes its cardinality $n+2$.
Thus \texttt{actual\_design} proves the full defining incidence
condition, not only a numerical parameter identity.

\texttt{IsAutomorphism} is the actual block-preservation biconditional
for a permutation of the point set. The theorem
\texttt{permutation\_preserves} proves it for every permutation.
Ordered triples of distinct points are represented by embeddings
\texttt{Fin 3} into the point set. The theorem
\texttt{three\_transitive} constructs a bijection of their ranges and
uses the finite-set extension theorem to produce an actual permutation
with the required action on all three entries.

\texttt{Isomorphic} requires an actual point bijection preserving
block membership in both directions. The theorem
\texttt{isomorphic\_iff} proves that two family members are isomorphic
exactly when their natural-number parameters are equal.
The formal quotient by this isomorphism relation is proved infinite
by an injective map from $\mathbb N$. Hence the final theorem
\texttt{conjecture\_00000008411} includes actual designs, actual
$3$-transitivity, the inequalities $3<k<v$, and an infinite type of
isomorphism classes.

\section{Verification}
Lean 4.19.0 is used with Mathlib commit
\texttt{c44e0c8ee63ca166450922a373c7409c5d26b00b}.
Run \texttt{lake build}, then \texttt{lake env lean Main.lean} in the
\texttt{lean} directory. The axiom audits of the design theorem,
transitivity theorem and final theorem contain only the standard logical
axioms \texttt{propext}, \texttt{Classical.choice} and \texttt{Quot.sound}.
There are no admitted proofs, extra axioms, native evaluation shortcuts,
or auxiliary computational dependencies.
\end{document}
